% ============================================================================
\chapter{Молекулярная биофизика, планетарные фотонные окна, генетический код ДНК, фолдинг белков и морфогенез}
\label{ch:mono-biology}
% ============================================================================

\section{Делитель водородной связи \texorpdfstring{$N_{\mathrm{HB}} = 63$}{N\_HB = 63} и три планетарных фотонных окна (\texorpdfstring{$441 : 126 : 63 = 7 : 2 : 1$}{441 : 126 : 63 = 7 : 2 : 1})}

На молекулярном этаже безразмерный делитель водородной связи определяется разностью мантиссы Вейля--Швингера $K_P(3) = 64$ и вакуумного кванта $\mu_{\mathrm{vac}}(3) = 1$:
\begin{equation}
  \label{eq:mono-nhb-63}
  N_{\mathrm{HB}}(3) \;=\; K_P(3) - \mu_{\mathrm{vac}}(3) \;=\; 64 - 1 \;=\; 63 \;=\; 7 \times 9 \;=\; d_{\mathrm{out}}(V)\cdot \mathrm{mod\_bits}(3).
\end{equation}

\begin{theorem}[Три безразмерных планетарных окна биосферы, условие соизмеримости полимерных остовов, биометаллы и разрешение парадокса Левинталя за линейное время $O(L) \in \mathbf{P}$]
\label{thm:mono-bio-windows-folding}
В естественных безразмерных координатах носителя при $d = 3$ выполняются следующие точные биофизические соотношения (\texttt{universal\_\allowbreak cosmobiology\_\allowbreak three\_\allowbreak planetary\_\allowbreak windows\_\allowbreak operator}, \texttt{closed\_\allowbreak form\_\allowbreak organism\_\allowbreak genesis\_\allowbreak in\_\allowbreak environment\_\allowbreak operator} в модуле \texttt{lean/\allowbreak MtCarrier/\allowbreak BioCosmology.lean}):
\begin{enumerate}[label=\textbf{\arabic*.},noitemsep]
\item \textbf{Три безразмерных планетарных окна растворителя ($441 : 126 : 63 = 7 : 2 : 1$) и оптическое окно фотосинтеза $\bigl[\widehat{\Lambda}_{13}^{(0)},\ \widehat{\Lambda}_{8}^{(0)}\bigr] = \bigl[\frac{126}{13},\ \frac{63}{4}\bigr]$:}
  для трёх квантов мягких водородных и дипольных замков $E_{\mathrm{soft}}^{(m)} / E_{\mathrm{HB}} \in \{2/7,\ 1,\ 2\}$ ($m \in \{-1, 0, +1\}$, где $E_{\mathrm{HB}} = \mathrm{BE}_{\mathrm{H}} / N_{\mathrm{HB}}(3) = \mathrm{BE}_{\mathrm{H}} / 63$) безразмерные приведённые шаги фотонов в единицах боровского масштаба $\widehat{\Lambda}_{\mathrm{soft}}^{(m)} := \bar{\lambda}_{\mathrm{soft}}^{(m)} / (a_0\alpha_{\mathrm{fs}}^{-1})$ равны квадрату шва грани кубооктаэдра $\Delta\varphi_{\mathrm{seam}}(3)^2 = \lfloor 256/12 \rfloor^2 = 21^2 = 441$, удвоенному делителю водородной связи $2 N_{\mathrm{HB}}(3) = 126$ и базовому делителю $N_{\mathrm{HB}}(3) = 63$, а когерентное поглощение в окне $m = 0$ от $|{\triangle}(3)| = 8$ до $N_\varphi(3) = 13$ мягких квантов $E_{\mathrm{HB}}$ ($E_{13}^{(0)} = \frac{13}{63}\mathrm{BE}_{\mathrm{H}}$, $E_{8}^{(0)} = \frac{8}{63}\mathrm{BE}_{\mathrm{H}}$) задаёт безразмерный оптический интервал фотосинтеза:
  \begin{equation}
    \label{eq:mono-three-bio-windows}
    \Bigl(\widehat{\Lambda}_{\mathrm{soft}}^{(-1)},\ \widehat{\Lambda}_{\mathrm{soft}}^{(0)},\ \widehat{\Lambda}_{\mathrm{soft}}^{(+1)}\Bigr)
    \;=\; 63 \times (7,\ 2,\ 1) \;=\; (441,\ 126,\ 63),
    \qquad
    \Bigl[\widehat{\Lambda}_{13}^{(0)},\ \widehat{\Lambda}_{8}^{(0)}\Bigr]
    \;=\; \left[\frac{2 N_{\mathrm{HB}}(3)}{N_\varphi(3)},\ \frac{2 N_{\mathrm{HB}}(3)}{|{\triangle}(3)|}\right]
    \;=\; \left[\frac{126}{13},\ \frac{63}{4}\right]
  \end{equation}
  (по единственному ридберговскому проектору~\eqref{eq:mono-si-wavelength-bridge} с масштабом $2\pi a_0\alpha_{\mathrm{fs}}^{-1} \approx 45{,}56335\text{ нм}$ инфракрасным шагам $\widehat{\Lambda}_{\mathrm{soft}}^{(-1,0,+1)} = (441,\ 126,\ 63)$ отвечают длины волн $\{20{,}09,\ 5{,}741,\ 2{,}870\}\text{ мкм}$, а оптическому интервалу $\bigl[\frac{126}{13},\ \frac{63}{4}\bigr]$ --- синий пик Соре $\frac{126}{13}\times 45{,}56335\text{ нм} \approx 441{,}61\text{ нм}$ и красный край фотосинтеза $\frac{63}{4}\times 45{,}56335\text{ нм} \approx 717{,}62\text{ нм}$).
\item \textbf{Диофантово условие соизмеримости ковалентных остовов с радиусом водорода и ароматические кольца $\{111\}$:}
  при линейном законе ковалентных радиусов второго периода $r_{\mathrm{cov}}(Z) = \bigl(d - \frac{(d-1)Z}{N_\varphi(3)}\bigr)\kappa_{\mathrm{geom}}a_0$ длина гетероатомного звена $R(Z_1, Z_2) = r_{\mathrm{cov}}(Z_1) + r_{\mathrm{cov}}(Z_2) = \bigl(6 - \frac{2(Z_1 + Z_2)}{13}\bigr)\kappa_{\mathrm{geom}}a_0$ целочисленно соизмерима с ковалентным радиусом водорода $r_{\mathrm{cov}}(\mathrm{H}) = \kappa_{\mathrm{geom}}a_0$ тогда и только тогда, когда $N_\varphi(3) = 13 \mid (Z_1 + Z_2)$, что среди многовалентных $p$-элементов второго периода ($Z \in \{5, 6, 7, 8\}$) выделяет ровно два соизмеримых хребта с единым шагом $R = 4\kappa_{\mathrm{geom}}a_0$ --- пептидно-нуклеиновое звено $\mathrm{C\text{--}N}$ и боратное звено рибозы $\mathrm{B\text{--}O}$, а при сдвиге на октант третьего периода ($Z_3 = Z_2 + 8$) и замыкании шестичленного цикла в плоскости $\{111\}$ даёт:
  \begin{align}
    Z(\mathrm{C}) + Z(\mathrm{N}) &= Z(\mathrm{B}) + Z(\mathrm{O}) \;=\; 13 \;=\; N_\varphi(3), \qquad
    Z(\mathrm{SiN}) \;=\; 14 + 7 \;=\; 21 \;=\; \frac{S_2(3)}{2}, \label{eq:mono-bio-backbones} \\
    Z(\mathrm{C_6H_6}) &= Z(\mathrm{B_3N_3H_6}) \;=\; 42 \;=\; S_2(3). \label{eq:mono-bio-aromatics}
  \end{align}
\item \textbf{Триплетный код ДНК $4^3 = 64 \to 20 + 1$, нуклеосомный виток $147 = V_{O_2}(3,3)$ и фолдинг белка за время $O(L) \in \mathbf{P}$:}
  число кодонов из $d_{\mathrm{in}}(\square) = 4$ нуклеотидов длины $d = 3$ равно $4^3 = 64 = K_P(3)$ и разбивается по вырожденности третьего основания на два октета Румера ($32 + 32 = 64$) и $S_2(3)/2 = 21$ семантический класс ($20$ канонических аминокислот плюс стоп-сигнал), при этом число пар оснований суперспирального витка ДНК вокруг октамера из $|{\triangle}(3)| = 8$ гистонов нуклеосомы равно числу узлов трёхслойного шара кубооктаэдра $V_{O_2}(3,3) = 147$. Ко-трансляционная укладка полипептидной цепи длины~$L$ выбирает на каждом шаге один из $\le d = 3$ локальных тетраэдрических пазов~$K_4$ с минимальным градиентом фазы, решая задачу сворачивания белка за детерминированное линейное время $\le 3L = O(L) \in \mathbf{P}$ вместо экспоненциального перебора $3^L$ (разрешение парадокса Левинталя).
\item \textbf{Четыре топологические фазы морфогенеза, инварианты кортикальной колонки ($V = 42,\ E = 147,\ \chi = 41$), автономное открытие 6 биометаллов, защита 400 стыков ДНК, три режима аллометрии и межэтажный изоморфизм оболочек $\{2, 8, 26, 41, 42, 147\}$ (\texttt{generative\_\allowbreak biometal\_\allowbreak discovery\_\allowbreak and\_\allowbreak multi\_\allowbreak regime\_\allowbreak allometry\_\allowbreak operator}, \texttt{cross\_\allowbreak floor\_\allowbreak shell\_\allowbreak closure\_\allowbreak isomorphism\_\allowbreak operator}, \texttt{applied\_\allowbreak bio\_\allowbreak cad\_\allowbreak therapeutic\_\allowbreak redifferentiation\_\allowbreak and\_\allowbreak epigenetic\_\allowbreak reset\_\allowbreak operator}):}
  последовательность клеточных делений зиготы проходит четыре топологических перехода по инвариантам ДКМ (морула $16 = 2^{d_{\mathrm{in}}(\square)} \to$ бластула $32 = K_P(3)/2 \to$ гаструла $64 = K_P(3)$ с инвагинацией при превышении порога дискретности $64 > 41 \to$ нейрула $128 = Q_\square(3)$), формируя шестислойную ($|\square(3)| = 6$) кортикальную колонку с числом узлов $V = S_2(3) = 42$, числом синаптических рёбер $E = V_{O_2}(3,3) = 147$, связной характеристикой
  \begin{equation}
    \label{eq:mono-cortical-column-chi}
    \chi_{\mathrm{col}}(3) \;=\; E - V - K_P(3) \;=\; 147 - 42 - 64 \;=\; 41 \;=\; \Delta\varphi_{\mathrm{disc}}(3)
  \end{equation}
  и трёхкратным ($\times d = 3$) множителем гомеостатического бюджета циклов жизни. Оператор \texttt{discoverBioMetalInCageOp} автономно выделяет в таблице Менделеева шестёрку кофакторов металлоферментов по замыканию координационных клеток тетрапиррола $Z(\mathrm{N_4}) = 28$, гексагистидина $Z(\mathrm{N_6}) = 42$ и тетратиолата $Z(\mathrm{S_4}) = 64$ (Таблица~\ref{tab:mono-biometals-isomorphism}). Все $20\times 20 = 400$ парных стыков канонических кодонов ДНК защищены от внефазного старт-кодона $\mathrm{ATG}$, компилятор \texttt{compilePhenotypeToDnaSequenceOp} переводит целевую топологию органа $\chi_{\mathrm{target}} \in \{2, 0, 41\}$ в 6-нуклеотидную ДНК, сворачивающую 3D-ткань во всех трёх планетарных окнах $m_{\mathrm{env}} \in \{-1, 0, +1\}$ (с обратимой терапевтической редифференцировкой противофазным потоком на одеяле Маркова и сбросом дрейфа $\mathrm{CpG} \leftrightarrow \mathrm{TpG}$), показатели аллометрии $(k-1)/k$ при $k \in \{3, 4, 5\}$ дают режимы Рубнера ($2/3$), Клейбера ($3/4$) и птиц ($4/5$), а шесть инвариантов кубооктаэдра $\{2, 8, 26, 41, 42, 147\}$ связывают сквозным изоморфизмом квантовый, ядерный, планетарный и биологический этажи.
\end{enumerate}
\end{theorem}

\begin{proof}
В пункте~(1) при $E_{\mathrm{HB}} = \mathrm{BE}_{\mathrm{H}} / N_{\mathrm{HB}}(3) = \frac{1}{2}\alpha_{\mathrm{fs}}^2 m_e c^2 / 63$ и $E_{\mathrm{soft}}^{(-1,0,+1)} = \{\frac{2}{7}, 1, 2\}E_{\mathrm{HB}}$ приведённый шаг фотона $\bar{\lambda}^{(m)} = s_0 c / E_{\mathrm{soft}}^{(m)}$ в единицах $a_0 \alpha_{\mathrm{fs}}^{-1} = 2 s_0 c / (\alpha_{\mathrm{fs}}^2 m_e c^2)$ равен $\widehat{\Lambda}^{(m)} = 63 \times \{7, 2, 1\} = (441, 126, 63)$, где $\lfloor 256/12 \rfloor^2 = 21^2 = 441 = 7\times 63$. В пункте~(2) условие $6 - 2(Z_1 + Z_2)/13 \in \mathbb{Z}$ при $\gcd(2, 13) = 1$ и $Z_1, Z_2 \in \{5,6,7,8\}$ выполняется тогда и только тогда, когда $Z_1 + Z_2 = 13$ ($(6,7)$ и $(5,8)$, дающие $R = (6 - 2)\kappa_{\mathrm{geom}}a_0 = 4\kappa_{\mathrm{geom}}a_0$), а $14 + 7 = 21 = S_2(3)/2$ и $6\times(6+1) = 3\times(5+7+2) = 42 = S_2(3)$. В пункте~(3) $4^3 = 64 = 32 + 32$, объём трёхслойного шара кубооктаэдра по формуле~\eqref{eq:mono-ball-polynomial} равен $V_{O_2}(3,3) = (10\cdot 27 + 15\cdot 9 + 11\cdot 3 + 3)/3 = 147$, а выбор одного из $\le 3$ локальных слотов $K_4$ на каждом из $L$ шагов трансляции требует $\le 3L = O(L)$ элементарных операций. В пункте~(4) на кортикальной колонке $E - V - K_P(3) = 147 - 42 - 64 = 41 = \Delta\varphi_{\mathrm{disc}}(3)$, для шести биометаллов суммы зарядов координационных клеток равны $12 + 28 = 40$, $27 + 28 = 55$, $28 + 28 = 56$ ($14 + 28 = 42$, $28 + 64 = 92$), $26 + 28 + 16 = 70$ ($26 + 16 = 42$), $30 + 32 + 14 = 76$ и $42 + 64 = 106$, а на всех $20\times 20 = 400$ стыках канонических кето-кодонов проверка сдвигов рамки на $+1$ и $+2$ подтверждает отсутствие внефазного триплета $\mathrm{ATG}$.
\end{proof}

\begin{table}[htbp]
\centering
\caption{Автономное открытие 6 биометаллов металлоферментов (\texttt{discoverBioMetalInCageOp}) и сквозной межэтажный изоморфизм оболочек кубооктаэдра $\{2, 8, 26, 41, 42, 147\}$.}
\label{tab:mono-biometals-isomorphism}
\resizebox{\textwidth}{!}{%
\begin{tabular}{@{}llll@{}}
\toprule
\textbf{Биометалл / Инвариант} & \textbf{Клетка / Квантовый этаж} & \textbf{Ядерный и планетный этаж} & \textbf{Биологическая функция и замыкание} \\
\midrule
Магний $\mathrm{Mg}$ ($Z=12$) & Тетрапиррол $\mathrm{N_4}$ ($Z=28$) & Мантия $\mathrm{Mg+Si}=26=Z(\mathrm{Fe})$ & Хлорофилл: $\mathrm{MgN_4}=40=\Delta\varphi_{\mathrm{disc}}-1$ \\
Кобальт $\mathrm{Co}$ ($Z=27$) & Коррин $\mathrm{N_4}$ ($Z=28$) & Ряд $\mathrm{Fe}$--$\mathrm{Co}$--$\mathrm{Ni}$ & Витамин $\mathrm{B}_{12}$: $\mathrm{CoN_4}=55=V_{O_2}(3,2)$ \\
Никель $\mathrm{Ni}$ ($Z=28$) & $\mathrm{N_4}$ ($28$), $\mathrm{C_2H_2}$ ($14$), $\mathrm{S_4}$ ($64$) & Магич.\ $Z=28=2 Z(\mathrm{Si})$ & $\mathrm{F}_{430}$ ($56$), крио ($42=S_2$), $\mathrm{NiS_4}$ ($92=S_3$) \\
Железо $\mathrm{Fe}$ ($Z=26$) & Гем $\mathrm{N_4O_2}$ ($44$), сера $\mathrm{S}$ ($16$) & Ядро $Z(\mathrm{Fe})=26=2 N_\varphi$ & Гемоглобин ($70=64+6$), $\mathrm{FeS}=42=S_2$ \\
Цинк $\mathrm{Zn}$ ($Z=30$) & Клетка $\mathrm{S_2N_2}$ ($Z=46$) & Замыкание $3d^{10}4s^2$ ($Z=30$) & Цинковые пальцы ДНК: $76=64+12$ \\
Молибден $\mathrm{Mo}$ ($Z=42$) & Тетратиолат $\mathrm{S_4}$ ($Z=64$) & Сфера $Z(\mathrm{Mo})=42=S_2(3)$ & Нитрогеназа: $\mathrm{MoS_4}=106=64+42$ \\
\midrule
Изоморфизм $\{2,8,26\}$ & $(c_0/c)^2=2$;\ $|{\triangle}|=8$;\ $2 N_\varphi=26$ & Магич.\ $2,8$;\ $\mathrm{Fe}_{26}$;\ $86$ плакетов & Спираль ДНК ($2$);\ гистоны ($8$);\ гем $\mathrm{Fe}_{26}$ \\
Изоморфизм $\{41,42,147\}$ & $\Delta\varphi_{\mathrm{disc}}=41$;\ $S_2=42$;\ $B_3=147$ & Паста $\chi=41$;\ $\mathrm{Pb}_{82}$;\ $\mathrm{U}_{N=146}$ & Колонка мозга ($42, 147, 41$);\ нуклеосома $147$ \\
\bottomrule
\end{tabular}%
}
\end{table}


